% Part: propositional-logic
% Chapter: syntax-and-semantics
% Section: introduction

\documentclass[../../../include/open-logic-section]{subfiles}

\begin{document}

\olfileid{pl}{syn}{int}

\olsection{مقدمة}

موضوع منطق القضايا !!{formula}s تُركَّب من
!!{propositional variable}s بالروابط القضوية
$\lnot$ و$\land$ و$\lor$ و$\lif$ و$\liff$. ومعنى
!!a{propositional variable}~${\OLId{latin-ordinary}{p}}$ في التصور الأول أنه يقوم مقام جملة أو قضية تقبل الصدق والكذب.
فإذا عُلمت قيمة الصدق لكل !!{propositional variable} في
!!a{formula}، عُلمت قيم الصدق لسائر !!{formula}s المؤلَّفة منها بهذه الروابط. لهذا قيل إن منطق القضايا
\emph{دالّ صدقيًا}: إذ تُفسَّر دلالته بدوال قيم الصدق. وليس من بحثنا هنا ما يزيد على مجرد الصدق والكذب؛ فلا نفرّق بين الصادق بالضرورة والصادق على سبيل العرض، ولا بين المعلوم صدقه وغيره، ولا بين الصادق الآن وما صدق من قبل أو سيصدق من بعد. وليس عندنا إلا قيمتان: الصدق
($\True$) والكذب ($\False$). فلا ندخل في هذا البحث عبارة لا تصدق ولا تكذب، ولا عبارة تصدق نصف صدق.
وكذلك لا نأخذ من الروابط إلا ما تتعيّن به قيمة صدق
!!a{formula} تعيّنًا تامًا من قيم صدق أجزائها، دون الرجوع إلى معانيها مثلًا. ومن ثم ينبغي التمييز: فكون الشرطيات الإنجليزية دالّةً صدقيًا بهذا المعنى موضع خلاف؛ وأما الشرط المادي~$\lif$ فدالّ صدقيًا. وللشرطيات التي لا تكون كذلك أنواع أخرى من المنطق.

أول ما نحتاج إليه لبناء نظرية منطق القضايا الدالّ صدقيًا وما وراء نظريته ضبطُ تركيب عباراته ودلالتها. نختار هنا طريقة لبناء
!!{formula}s من !!{propositional variable}s بالروابط، وليست هي الطريقة الوحيدة. فقد تختلف الأنظمة في رموزها، وفي الروابط التي تُتخذ أولية، وفي استعمال الأقواس؛ بل قد تستغني عنها، كما في التدوين المسمّى بولنديًا.
ومهما اختلفت هذه الاختيارات، فالمشترك أن قواعد التكوين تحدد مجموعة
!!{formula}s \emph{استقرائيًا}. وإذا أُحكم وضعها، لم يكن لكل تعبير في أصله إلا وجه واحد من التوليد بحسبها. وهذا هو كون التعابير
\emph{ذات قراءة وحيدة}. وبه يتأتى أن نعرّف معانيها بالطريق نفسه، أعني التعريف الاستقرائي.

وأما الدلالة فشأنها تعيين معاني التعابير، وأصلها هنا الإرضاء في
!!a{valuation}. ذلك أن !!^a{valuation}~$\pAssign{{\OLId{latin-ordinary}{v}}}$ يسند
$\True$ و$\False$ إلى !!{propositional variable}s؛ فيعيّن كل
!!{valuation} قيمة الصدق $\pValue{{\OLId{latin-ordinary}{v}}}(!A)$ لكل
!!{formula}~$!A$. وتكون !!^a{formula} مُرضاة في
!!a{valuation}~$\pAssign{{\OLId{latin-ordinary}{v}}}$ متى وفقط متى كانت
$\pValue{{\OLId{latin-ordinary}{v}}}(!A) = \True$، ونرمز إلى ذلك بـ
$\pSat{{\OLId{latin-ordinary}{v}}}{!A}$. ولنا أن نعرّف هذه العلاقة بالاستقراء على بنية~$!A$؛ فنستعمل دوال الروابط الصدقية لبيان إرضاء
$!A \land !B$، مثلًا، انطلاقًا من إرضاء $!A$ و~$!B$ أو عدم إرضائهما.

ومن علاقة الإرضاء $\pSat{{\OLId{latin-ordinary}{v}}}{!A}$ للجمل تتفرع المعاني الدلالية الأساسية: التحصيل الحاصل والاستلزام وقابلية الإرضاء. فإن أرضى كل
!!{valuation} ال!!^a{formula}، فهي تحصيل حاصل، ونكتب
$\Entails !A$؛ ومعناه $\pValue{{\OLId{latin-ordinary}{v}}}(!A) = \True$ لكل~$\pAssign{{\OLId{latin-ordinary}{v}}}$.
وإذا كانت مجموعة من !!{formula}s تستلزمها، كتبنا
${\OLId{greek-variable}{Gamma}} \Entails !A$؛ ومعناه أن كل !!{valuation} يرضي جميع
!!{formula}s في~${\OLId{greek-variable}{Gamma}}$ يرضي~$!A$ كذلك.
وأما قابلية إرضاء مجموعة من !!{formula}s فأن يوجد
!!{valuation} واحد يرضي جميع !!{formula}s فيها معًا.
ولما كان تعريف !!{formula}s استقرائيًا، وكان تعريف الإرضاء نفسه بالاستقراء على بنية
!!{formula}s، أمكن أن نبرهن بالاستقراء خواص هذه الدلالة والروابط بين مفاهيمها.


\end{document}
